\section*{Sets and Indices}

\begin{itemize}
    \item Gasoline products $g \in G = \{\mathrm{I}, \mathrm{II}\}$.
    \item Purchase tiers for crude A $k \in K = \{1,2,3\}$.
\end{itemize}

\section*{Parameters}

\begin{itemize}
    \item Minimum required proportion of crude A in gasoline $g$:
    \[
    \alpha_{\mathrm{I}} \text{ (code: \texttt{min\_A\_I})}, \qquad
    \alpha_{\mathrm{II}} \text{ (code: \texttt{min\_A\_II})}.
    \]
    \item Selling price of gasoline $g$:
    \[
    r_{\mathrm{I}} \text{ (code: \texttt{price\_I})}, \qquad
    r_{\mathrm{II}} \text{ (code: \texttt{price\_II})}.
    \]
    \item Initial inventory of crude oils:
    \[
    I_A \text{ (code: \texttt{inv\_A})}, \qquad
    I_B \text{ (code: \texttt{inv\_B})}.
    \]
    \item Maximum total purchasable amount of crude A:
    \[
    \bar P_A \text{ (code: \texttt{max\_purchase\_A})}.
    \]
    \item Unit purchase cost of crude A in tier $k$:
    \[
    c_1 \text{ (code: \texttt{cost\_tier1})}, \quad
    c_2 \text{ (code: \texttt{cost\_tier2})}, \quad
    c_3 \text{ (code: \texttt{cost\_tier3})}.
    \]
    \item Tier capacity (implemented identically for all three tiers):
    \[
    U = 500.
    \]
\end{itemize}

\section*{Decision Variables}

\begin{itemize}
    \item Amount of crude A blended into gasoline I and II:
    \[
    a_1 \ge 0 \text{ (code: \texttt{a1})}, \qquad
    a_2 \ge 0 \text{ (code: \texttt{a2})}.
    \]
    \item Amount of crude B blended into gasoline I and II:
    \[
    b_1 \ge 0 \text{ (code: \texttt{b1})}, \qquad
    b_2 \ge 0 \text{ (code: \texttt{b2})}.
    \]
    \item Purchased amount of crude A in each tier:
    \[
    p_1 \in [0,U] \text{ (code: \texttt{p1})}, \qquad
    p_2 \in [0,U] \text{ (code: \texttt{p2})}, \qquad
    p_3 \in [0,U] \text{ (code: \texttt{p3})}.
    \]
    \item Binary variables enforcing sequential use of purchase tiers:
    \[
    y_1 \in \{0,1\} \text{ (code: \texttt{y1})}, \qquad
    y_2 \in \{0,1\} \text{ (code: \texttt{y2})}.
    \]
\end{itemize}

\section*{Objective}

Maximize total gasoline revenue minus crude A purchase cost:
\[
\max \;
r_{\mathrm{I}}(a_1+b_1) + r_{\mathrm{II}}(a_2+b_2)
- c_1 p_1 - c_2 p_2 - c_3 p_3.
\]

\section*{Constraints}

Tier-ordering constraints for crude A purchases:
\begin{align}
p_1 &\le U, \\
p_2 &\le U y_1, \\
p_3 &\le U y_2, \\
p_1 &\ge U y_1, \\
p_2 &\ge U y_2.
\end{align}

Crude availability constraints:
\begin{align}
a_1 + a_2 &\le I_A + p_1 + p_2 + p_3, \\
b_1 + b_2 &\le I_B.
\end{align}

Minimum crude-A proportion requirements:
\begin{align}
a_1 &\ge \alpha_{\mathrm{I}}(a_1+b_1), \\
a_2 &\ge \alpha_{\mathrm{II}}(a_2+b_2).
\end{align}

Variable domains:
\begin{align}
a_1,a_2,b_1,b_2 &\ge 0, \\
0 \le p_1,p_2,p_3 &\le U, \\
y_1,y_2 &\in \{0,1\}.
\end{align}

\section*{Notes on Implementation}

\begin{itemize}
    \item The implemented model is a mixed-integer linear program solved with Gurobi through PuLP.
    \item The tiered-purchase logic is modeled exactly as in the code: tier 2 can be used only if tier 1 is full, and tier 3 can be used only if tier 2 is full. This is enforced by
    \[
    p_1 \ge U y_1,\quad p_2 \le U y_1,\quad
    p_2 \ge U y_2,\quad p_3 \le U y_2.
    \]
    Thus, if $y_1=1$ then necessarily $p_1=U$, and if $y_2=1$ then necessarily $p_2=U$.
    \item Although parameter $\bar P_A$ (code: \texttt{max\_purchase\_A}) is loaded from data, it is not used in any constraint. The implemented upper limit on total purchased crude A is instead implied by the three tier bounds, i.e.,
    \[
    p_1+p_2+p_3 \le 3U = 1500.
    \]
    \item No demand limits, processing capacity limits, or disposal variables are included. All available crude B may be used if profitable, and crude inventories need not be exhausted.
\end{itemize}
